% ============================================================
% 中文摘要
% ============================================================

随着电子信息产业的快速发展，电子元器件选型与供应链风险管理已成为制约产品研发效率和可靠性的关键环节。传统的人工选型流程依赖工程师经验，耗时长、易出错，且缺乏对供应链风险的系统性评估。本文面向eZ-PLM电子元器件管理平台，设计并实现了一套基于大语言模型（LLM）与检索增强生成（RAG）技术的智能元器件选型与风险评估Agent系统。

该系统采用"自然语言需求解析—多维知识检索—智能评分排序—风险识别—报告生成"五阶段Pipeline架构，并在此基础上引入LangChain ReAct Agent范式实现多工具自主编排与多轮对话。系统核心创新包括：（1）提出LLM优先解析+规则兜底的双重需求理解机制，支持中英文自然语言输入，覆盖Buck/Boost/LDO三类电源拓扑；（2）构建基于ChromaDB的工程知识向量库，通过语义检索将设计公式、标准规范等专业知识注入选型决策链，实现知识增强型推理；（3）设计五维加权评分模型（参数匹配、供应链稳定性、应用场景适配度、设计成熟度、证据置信度），支持纯规则与LLM增强双模式切换；（4）定义PartIR/RiskIR/TopologyIR三套标准化中间表示（IR），实现与eZ-PLM平台的数据互通。

实验结果表明，系统在10个DC-DC降压选型测试用例上实现100\%的通过率。RAG知识增强使拓扑分析报告的工程规范性显著提升，能够自动输出包含电感计算公式、热性能估算和PCB Layout指导的完整设计文档。ReAct Agent模式支持国产替代追问等多轮交互场景，单轮工具调用控制在2-4次以内。系统生成的BOM清单、风险评估报告和拓扑分析文档可直接用于工程评审，验证了AI Agent在电子工程领域的实用价值。
